\documentclass[conference,compsoc]{IEEEtran}
\usepackage{listings}
\usepackage{xcolor}
\usepackage{graphicx}
\usepackage{subfig}
\usepackage{amsmath}

\definecolor{codegreen}{rgb}{0,0.6,0}
\definecolor{codegray}{rgb}{0.5,0.5,0.5}
\definecolor{codepurple}{rgb}{0.58,0,0.82}
\definecolor{backcolour}{rgb}{0.95,0.95,0.92}

\lstdefinestyle{mystyle}{
    backgroundcolor=\color{backcolour},   
    commentstyle=\color{codegreen},
    keywordstyle=\color{magenta},
    numberstyle=\tiny\color{codegray},
    stringstyle=\color{codepurple},
    basicstyle=\ttfamily\footnotesize,
    breakatwhitespace=false,         
    breaklines=true,                 
    captionpos=b,                    
    keepspaces=true,                 
    numbers=left,                    
    numbersep=5pt,                  
    showspaces=false,                
    showstringspaces=false,
    showtabs=false,                  
    tabsize=2
}

\lstset{style=mystyle}

% *** CITATION PACKAGES ***
%
\ifCLASSOPTIONcompsoc
  % IEEE Computer Society needs nocompress option
  % requires cite.sty v4.0 or later (November 2003)
  \usepackage[nocompress]{cite}
\else
  % normal IEEE
  \usepackage{cite}
\fi

% correct bad hyphenation here
\hyphenation{op-tical net-works semi-conduc-tor}


\begin{document}
%
% Title
\title{Flow field through a porous media\\
GROUP NAME AND NUMBER}


% names and affiliations
\author{\IEEEauthorblockN{Name surname}
\IEEEauthorblockA{Student id\\dept./bsc course\\
email}
\and
\IEEEauthorblockN{Name surname}
\IEEEauthorblockA{Student id\\dept./bsc course\\
email}
\and
\IEEEauthorblockN{Name surname}
\IEEEauthorblockA{Student id\\dept./bsc course\\
email}}


% make the title area
\maketitle
\thispagestyle{plain}
\pagestyle{plain}

% As a general rule, do not put math, special symbols or citations
% in the abstract
\begin{abstract}
Optional. Short summary of the work - Max 120 words.
\end{abstract}

% no keywords




\section{Introduction}
% no \IEEEPARstart
Brief explanation of the overall assignment, of the performance reached, and of the components done.

\section{Part 1 - Theoretical aspects}

\subsection*{A - Content aligned with point A of the rubric/assignment}

Example of inline formula: $y = f(x)$. A block formula instead is as follows
\begin{equation}
    y = f(x).
\end{equation}

\subsection*{B - Content aligned with point B of the rubric/assignment}

Example of figure, which is referenced as Figure~\ref{fig:example-label}.

\begin{figure}[h]
    \centering
    \includegraphics[width=1\linewidth]{fig/example_image.png}
    \caption{Caption example}
    \label{fig:example-label}
\end{figure}


Example of a subfigure, which elements are referenced as Figure~\ref{fig:first_case} and \ref{fig:full_sub_fig}

\begin{figure}[h]
    \centering
    \subfloat[Subcaption example]{\includegraphics[width=0.49\linewidth]{fig/example_subfig_1.png}%
    \label{fig:first_case}}
    \hfil
    \subfloat[Subcaption example]{\includegraphics[width=0.49\linewidth]{fig/example_subfig_2.png}%
    \label{fig:second_case}}
    \caption{Full caption example.}
    \label{fig:full_sub_fig}
\end{figure}

Alternatively display the same figures spanning across the two text columns as in Figure~\ref{fig:full_sub_fig_alt}

\begin{figure*}[h]
    \centering
    \subfloat[Subcaption example]{\includegraphics[width=0.49\linewidth]{fig/example_subfig_1.png}%
    \label{fig:first_case_alt}}
    \hfil
    \subfloat[Subcaption example]{\includegraphics[width=0.49\linewidth]{fig/example_subfig_2.png}%
    \label{fig:second_case_alt}}
    \caption{Full caption example.}
    \label{fig:full_sub_fig_alt}
\end{figure*}

Bibliography references use instead the \texttt{\\cite} keyword, e.g. \cite{IEEEhowto:kopka}.

\section{Part 2 - Implementation and code}

Discuss key design choices... include key code snippets with \texttt{lstlisting}.
\begin{lstlisting}[language=Python]
import numpy as np
    
def incmatrix(genl1,genl2):
    m = len(genl1)
    n = len(genl2)
    M = None #to become the incidence matrix
    VT = np.zeros((n*m,1), int)  #dummy variable
    
    #compute the bitwise xor matrix
    M1 = bitxormatrix(genl1)
    M2 = np.triu(bitxormatrix(genl2),1) 

    for i in range(m-1):
        for j in range(i+1, m):
            [r,c] = np.where(M2 == M1[i,j])
            for k in range(len(r)):
                VT[(i)*n + r[k]] = 1;
                VT[(i)*n + c[k]] = 1;
                VT[(j)*n + r[k]] = 1;
                VT[(j)*n + c[k]] = 1;
                
                if M is None:
                    M = np.copy(VT)
                else:
                    M = np.concatenate((M, VT), 1)
                
                VT = np.zeros((n*m,1), int)
    
    return M
\end{lstlisting}

Showcase convergence/absence of overfit...





\section{Conclusion}
Always sum up your message and results in a concluding section.


% references section

% can use a bibliography generated by BibTeX as a .bbl file
% BibTeX documentation can be easily obtained at:
% http://mirror.ctan.org/biblio/bibtex/contrib/doc/
% The IEEEtran BibTeX style support page is at:
% http://www.michaelshell.org/tex/ieeetran/bibtex/
%\bibliographystyle{IEEEtran}
% argument is your BibTeX string definitions and bibliography database(s)
%\bibliography{IEEEabrv,../bib/paper}
%
% <OR> manually copy in the resultant .bbl file
% set second argument of \begin to the number of references
% (used to reserve space for the reference number labels box)
\begin{thebibliography}{1}

\bibitem{IEEEhowto:kopka}
H.~Kopka and P.~W. Daly, \emph{A Guide to \LaTeX}, 3rd~ed.\hskip 1em plus
  0.5em minus 0.4em\relax Harlow, England: Addison-Wesley, 1999.

\end{thebibliography}

\begin{appendices}
\section{Appendix A title}
The contents...

\end{appendices}



% that's all folks
\end{document}


